% Pursuing Stacks (cotes 134-2 to 134-8): archive cote 134-5, site page 14 (PDF page 15), typescript page 201.
% typescript. Draft compared: ps4.tex:612-670.
% AI-transcribed by Claude Opus 5.5 (claude-opus-5-5, adaptive thinking, effort high), Message Batches API, batch ps2 (msgbatch_01NVwr1Yuo8qJf5rykeGBXPK), 2026-11-12; not yet checked by eye.
\scanpage{134-5/14}
\mspage{201}
main ingredients of a (\inserted{by then} still very misty) notion of a triangulated category
of sorts\hand{)}. This feeling arose from my  ponderings on stratified structures,
and the ``screwing together'' of such structures in terms of simple building
blocks (essentially,  various types of ``tubes'' associated to such structu\xout{r}%
res,  related to each other by various \inserted[typed]{proper} maps which are either inclusions,
or  -  in the equisingular case \inserted[typed]{at any rate} -  fiber maps).  This\hand{``}screwing together
\inserted{operation''} could be expressed as being a direct limit of a certain finite system of
spaces.   In the cases I was most interested in (namely the Mumford-Deligne
compactifications $M^{\wedge}_{g,\hand{\nu}}$  of the modular multiplicities  $M_{g,\hand{\nu}}$),  these
spaces or\hand{``}tubes\hand{''} have exceedingly simple homotopy types  -  they \xout{wwww} \inserted[typed]{are} just
$K(\hand{\pi},1)$-spaces,  where \xout{ex} \inserted[typed]{each} \hand{$\pi$} is a Teichmüller-type discrete group (practical-%
ly,  a product of usual Teichmüller groups).  It then occured to me that
the \inserted[typed]{whole} homotopy type of $M^{\wedge}_{g,\hand{\nu}}$ ,  or of any locally closed union of strata\hand{,}
or (more generally still) of ``the'' tubular neighbourhood of \xout{the} such a
union in any larger one\hand{,} \inserted[typed]{etc} -  that all these homotopy types should be ex-%
pressible in terms of the given \xout{direct} system of spaces,  and more accura-%
tely still,  just in terms of the corresponding system of fundamental
groupoids (embodying their homotopy types).  In this situation,  what I
was mainly out for,  was precisely an accurate and workeable description
of this direct system of groupoids  (  which could be viewed as just one
section of the whole ``Teichmüller tower'' of Teichmüller groupoïds ...).
Thus,  it was a rewarding extra feature of the situation (by then just an
expectation\hand{,} as a matter of fact),  that such a description should at the
same time yie\over{\unsure{l}}{l}d a ``purely algebraic'' description of the homotopy types
of all the spaces (rather,  multiplicities,to be wholly accurate) which
I could think of in terms of the natural stratification of $M^{\wedge}_{g,\hand{\nu}}$ .  The\hand{re}
was an awareness that this operation on homotopy types could not be
described simply in terms of a functor  I \hand{$\longrightarrow$} (Hot) ,  where I is the in-%
dexing category,  that a functor \inserted{$i\mapsto X_i$ :} I \hand{$\longrightarrow$} M  (where M is some model category
such as (Spaces)  or (Cat) ) should be available in order to define an
``integra\over{\unsure{l}}{t}ed homotopy type''  \hand{$\int$}$X_i$ .  This justified feeling got somewhat
blurred lately,  for a little while,  by the definitely unreasonable ex-%
pectation that finite limits should exist in (Hot) after all,  why not !
It's enough to have a look though (which probably I did years ago and then
forgot in the meanwhile) to ma\over{\unsure{k}}{k}e sure they dont ...

\quad Whether or not this notion of ``integration of homotopy types'' is
more or less well known already under some name or other,  is'nt quite
clear to me.  It is'nt familiar to Ronnie Brown visibly,  but it seems he
heard about such a kind of thing,  without \inserted{his} being specific about it.  It
\margin{14}
